\chapter{โครงการลดก๊าซเรือนกระจกภาคสมัครใจ T-VER และกลไกคาร์บอนเครดิต}
\label{chap:ch08}

\section*{แผนบริหารการสอนประจำบทที่ 8}
\addcontentsline{toc}{section}{แผนบริหารการสอนประจำบทที่ 8}

\noindent\textbf{หัวข้อเนื้อหาประจำบท}
\begin{enumerate}
    \item กรอบความตกลงปารีสและกลไกคาร์บอนเครดิตในประเทศไทย
    \item ระเบียบวิธี T-VER สาขาพลังงานหมุนเวียนและชีวมวล
    \item ระเบียบวิธี T-VER สาขาป่าไม้และการเกษตร
    \item กระบวนการขอรับรองและการตรวจสอบความใช้ได้ของโครงการ
\end{enumerate}

\noindent\textbf{วัตถุประสงค์เชิงพฤติกรรม}
\begin{enumerate}
    \item อธิบายหลักการ ทฤษฎีฟิสิกส์ และแบบจำลองทางสิ่งแวดล้อมประจำบทที่ 8 ได้อย่างถูกต้อง
    \item วิเคราะห์ สังเคราะห์ และประยุกต์ใช้เครื่องมือตรวจวัดหรือกระบวนการแปรรูปพลังงานได้อย่างมีประสิทธิภาพ
    \item คำนวณ ประเมินผลทางเทอร์โมไดนามิกส์ ทัศนศาสตร์ หรือการลดก๊าซเรือนกระจกได้อย่างถูกต้องตามหลักวิชาการ
\end{enumerate}

\noindent\textbf{การวัดและประเมินผล}
\begin{table}[H]
\centering\small
\begin{tabularx}{\textwidth}{p{0.55\textwidth} p{0.25\textwidth} c}
\toprule
\textbf{เกณฑ์การประเมินผล} & \textbf{เครื่องมือวัดผล} & \textbf{สัดส่วนคะแนน} \\
\midrule
1. ทักษะการปฏิบัติการทดลองและการคำนวณ & แบบประเมินทักษะห้องแล็บ & 40\% \\
2. รายงานผลการทดลองและการวิเคราะห์ & การตรวจดาต้าชีทและรายงาน & 30\% \\
3. แบบฝึกหัดทบทวนและคำถามท้ายบท & แบบทดสอบท้ายบทเรียน & 20\% \\
4. จิตพิสัยและการมีส่วนร่วม & การเข้าเรียนและการตรงต่อเวลา & 10\% \\
\midrule
\textbf{รวมทั้งสิ้น} & & \textbf{100\%} \\
\bottomrule
\end{tabularx}
\end{table}
\newpage

% ==========================================================
% เนื้อหาบทเรียนประจำบทที่ 8
% ==========================================================

\section{กรอบความตกลงปารีสและกลไกคาร์บอนเครดิตในประเทศไทย}
\index{กรอบความตกลงปารีสและกลไกคาร์บอนเครดิตในประเทศไทย}

การศึกษาและทำความเข้าใจเกี่ยวกับกรอบความตกลงปารีสและกลไกคาร์บอนเครดิตในประเทศไทย (Paris Agreement and Thailand Carbon Credit Mechanism) มีความสำคัญอย่างยิ่งในงานฟิสิกส์สิ่งแวดล้อมและพลังงาน โดยมีรายละเอียดสาระสำคัญดังนี้

ภายใต้ความตกลงปารีส (Paris Agreement) มาตรา 6 ประเทศไทยได้ประกาศเป้าหมายความเป็นกลางทางคาร์บอน (Carbon Neutrality) ภายในปี ค.ศ. 2050 และการปล่อยก๊าซเรือนกระจกสุทธิเป็นศูนย์ (Net Zero GHG Emissions) ภายในปี ค.ศ. 2065

องค์การบริหารจัดการก๊าซเรือนกระจก (องค์การมหาชน) หรือ อบก. (TGO) ได้พัฒนากลไก โครงการลดก๊าซเรือนกระจกภาคสมัครใจตามมาตรฐานของประเทศไทย (Thailand Voluntary Emission Reduction Program; T-VER) เพื่อส่งเสริมให้ทุกภาคส่วนดำเนินกิจกรรมลดการปล่อยก๊าซเรือนกระจก และสามารถนำปริมาณก๊าซเรือนกระจกที่ลดหรือกักเก็บได้ ที่เรียกว่า คาร์บอนเครดิต (Carbon Credits มีหน่วยเป็นตันคาร์บอนไดออกไซด์เทียบเท่า; $tCO_2e$) ไปซื้อขายแลกเปลี่ยนเพื่อชดเชยการปล่อยก๊าซเรือนกระจก (Carbon Offsetting) ได้

\begin{casebox}[มาตรฐานการลดก๊าซเรือนกระจก T-VER ประเทศไทย]
โครงการลดก๊าซเรือนกระจกภาคสมัครใจต้องผ่านการตรวจสอบความใช้ได้และการทวนสอบโดยหน่วยงานอิสระ (VVB) ตามระเบียบวิธีของ อบก. เพื่อให้คาร์บอนเครดิตมีความโปร่งใสและน่าเชื่อถือในระดับสากล
\end{casebox}

\section{ระเบียบวิธี T-VER สาขาพลังงานหมุนเวียนและชีวมวล}
\index{ระเบียบวิธี T-VER สาขาพลังงานหมุนเวียนและชีวมวล}

การศึกษาและทำความเข้าใจเกี่ยวกับระเบียบวิธี T-VER สาขาพลังงานหมุนเวียนและชีวมวล (T-VER Methodology for Renewable Biomass Energy) มีความสำคัญอย่างยิ่งในงานฟิสิกส์สิ่งแวดล้อมและพลังงาน โดยมีรายละเอียดสาระสำคัญดังนี้

การพัฒนาโครงการผลิตพลังงานความร้อนจากชีวมวลและถ่านอัดแท่งเพื่อขอรับรองคาร์บอนเครดิต ต้องปฏิบัติตามระเบียบวิธี T-VER-METH-01 (การผลิตพลังงานความร้อนจากชีวมวลทดแทนเชื้อเพลิงฟอสซิล)

สมการคำนวณการลดการปล่อยก๊าซเรือนกระจก:
\[ ER\_y = BE\_y - PE\_y - LE\_y \]
โดยที่:
- $ER_y$ คือปริมาณการลดการปล่อยก๊าซเรือนกระจกในปี $y$ ($tCO_2e/year$)
- $BE_y$ คือการปล่อยก๊าซเรือนกระจกในกรณีฐาน (Baseline Emissions) จากการเผาไหม้เชื้อเพลิงฟอสซิลเดิม (เช่น น้ำมันเตา หรือถ่านหิน) เพื่อผลิตพลังงานความร้อนในปริมาณเท่ากัน:
\[ BE\_y = \frac{EG\_y}{\eta\_{\text{baseline}}} \times EF\_{\text{fossil}} \]
- $PE_y$ คือการปล่อยก๊าซเรือนกระจกจากการดำเนินโครงการ (Project Emissions) เช่น การใช้ไฟฟ้าในการเดินเครื่องจักรและเครื่องอัดแท่ง
- $LE_y$ คือการปล่อยก๊าซเรือนกระจกนอกขอบเขตโครงการ (Leakage Emissions) เช่น การขนส่งชีวมวลระยะไกล

\begin{flamebox}[นวัตกรรมพลังงานความร้อน มรภ.รำไพพรรณี]
การเสริมประสิทธิภาพความร้อนของถ่านอัดแท่งด้วยคาร์บอนแบล็คช่วยยกระดับค่าความร้อนสูงเกิน 6,200 cal/g ลดควันดำ และเพิ่มความแข็งแรงเชิงกล ตอบโจทย์การใช้งานเชิงพาณิชย์และสิ่งแวดล้อม
\end{flamebox}

\section{ระเบียบวิธี T-VER สาขาป่าไม้และการเกษตร}
\index{ระเบียบวิธี T-VER สาขาป่าไม้และการเกษตร}

การศึกษาและทำความเข้าใจเกี่ยวกับระเบียบวิธี T-VER สาขาป่าไม้และการเกษตร (T-VER Forestry and Agricultural Methodologies) มีความสำคัญอย่างยิ่งในงานฟิสิกส์สิ่งแวดล้อมและพลังงาน โดยมีรายละเอียดสาระสำคัญดังนี้

โครงการด้านการเกษตรและป่าไม้ (Forestry and Agriculture) มุ่งเน้นการเพิ่มพูนการกักเก็บก๊าซเรือนกระจกในชีวมวลพืชและดิน:
- โครงการปลูกป่าและอนุรักษ์ป่าไม้: คำนวณมวลชีวภาพเหนือพื้นดิน (Above-Ground Biomass; AGB) และชีวมวลใต้ดิน (Below-Ground Biomass; BGB) ผ่านสมการแอลโลเมตรี (Allometric equations)
- โครงการสวนไม้ผลและไม้ยืนต้นเศรษฐกิจ: ต้นทุเรียน ยางพารา และปาล์มน้ำมัน จัดเป็นไม้ยืนต้นที่มีอัตราการสังเคราะห์แสงและกักเก็บคาร์บอนในเนื้อไม้อย่างต่อเนื่อง
- โครงการฝังกลบถ่านชีวภาพ (Biochar Soil Amendment): ได้รับการรับรองภายใต้มาตรฐาน T-VER ขั้นสูง (Premium T-VER) ในฐานะเทคโนโลยีการกำจัดคาร์บอนไดออกไซด์ถาวร (Carbon Dioxide Removal; CDR)

\begin{casebox}[มาตรฐานการลดก๊าซเรือนกระจก T-VER ประเทศไทย]
โครงการลดก๊าซเรือนกระจกภาคสมัครใจต้องผ่านการตรวจสอบความใช้ได้และการทวนสอบโดยหน่วยงานอิสระ (VVB) ตามระเบียบวิธีของ อบก. เพื่อให้คาร์บอนเครดิตมีความโปร่งใสและน่าเชื่อถือในระดับสากล
\end{casebox}

\section{กระบวนการขอรับรองและการตรวจสอบความใช้ได้ของโครงการ}
\index{กระบวนการขอรับรองและการตรวจสอบความใช้ได้ของโครงการ}

การศึกษาและทำความเข้าใจเกี่ยวกับกระบวนการขอรับรองและการตรวจสอบความใช้ได้ของโครงการ (Validation, Verification and Credit Issuance) มีความสำคัญอย่างยิ่งในงานฟิสิกส์สิ่งแวดล้อมและพลังงาน โดยมีรายละเอียดสาระสำคัญดังนี้

ขั้นตอนการพัฒนาโครงการ T-VER ตามมาตรฐาน อบก.:
1. จัดทำเอกสารข้อเสนอโครงการ (Project Design Document; PDD) และกำหนดเส้นฐานการปล่อยก๊าซเรือนกระจก
2. การตรวจสอบความใช้ได้ (Validation) โดยผู้ประเมินภายนอกที่ขึ้นทะเบียน (Validation and Verification Body; VVB)
3. ยื่นขอขึ้นทะเบียนโครงการ (Project Registration) ต่อคณะกรรมการ อบก.
4. ดำเนินโครงการและติดตามผลการตรวจวัดจริง (Monitoring) ตามแผนที่กำหนดใน PDD
5. การทวนสอบผลการลดการปล่อยก๊าซ (Verification) โดย VVB ประจำปี
6. อบก. ออกใบรับรองคาร์บอนเครดิต (Credit Issuance) เข้าสู่บัญชีทะเบียน T-VER Registry พร้อมสำหรับการซื้อขายในตลาดคาร์บอน

\section{ขั้นตอนและวิธีปฏิบัติการทดลอง}
\index{ขั้นตอนการทดลองประจำบทที่ 8}

\subsection{การคำนวณคาร์บอนเครดิตของโครงการเปลี่ยนเชื้อเพลิงชีวมวล}
\begin{conceptbox}[การคำนวณคาร์บอนเครดิตของโครงการเปลี่ยนเชื้อเพลิงชีวมวล]
1. รวบรวมข้อมูลปริมาณการใช้เชื้อเพลิงชีวมวลถ่านอัดแท่งในโรงงานอุตสาหกรรม: 5,000 ตัน/ปี
2. วัดค่าความร้อนเฉลี่ยของถ่านอัดแท่ง: $22.0$ GJ/ตัน คำนวณพลังงานความร้อนรวม: $EG = 110,000$ GJ/ปี
3. เชื้อเพลิงฟอสซิลเดิมคือน้ำมันเตา (Fuel Oil Class A) มีค่าสัมประสิทธิ์การปล่อยก๊าซ $EF = 0.0774$ $tCO_2e/GJ$
4. ประสิทธิภาพหม้อไอน้ำเดิม $\eta = 80\%$ คำนวณการปล่อยกรณีฐาน: $BE = (110,000 / 0.80) \times 0.0774 = 10,642.5\text{ tCO}_2\text{e/ปี}$
5. การปล่อยจากการผลิตและขนส่งถ่านอัดแท่ง: $PE = 145.0$ $\text{tCO}_2\text{e/ปี}$, การรั่วไหล: $LE = 82.5$ $\text{tCO}_2\text{e/ปี}$
6. คำนวณปริมาณคาร์บอนเครดิตสุทธิ: $ER = 10,642.5 - 145.0 - 82.5 = 10,415.0\text{ tCO}_2\text{e/ปี}$
\end{conceptbox}

\subsection{การสำรวจและคำนวณชีวมวลและการกักเก็บคาร์บอนในสวนไม้ยืนต้น}
\begin{conceptbox}[การสำรวจและคำนวณชีวมวลและการกักเก็บคาร์บอนในสวนไม้ยืนต้น]
1. วางแปลงตัวอย่างถาวรขนาด $40 \times 40$ เมตร (0.1 ไร่ หรือ 1,600 ตร.ม.) ในสวนไม้ผล
2. วัดเส้นผ่านศูนย์กลางลำต้นที่ระดับอก (Diameter at Breast Height; DBH ที่ 1.3 เมตร) และความสูงต้น ($H$)
3. คำนวณมวลชีวภาพเหนือพื้นดิน (AGB, kg) ตามสมการแอลโลเมตรี: $AGB = 0.0632 \times (\text{DBH}^2 \cdot H)^{0.912}$
4. คำนวณมวลชีวภาพรวม: $TB = AGB \times (1 + R)$ โดยกำหนดอัตราส่วนราก $R = 0.20$
5. คำนวณปริมาณการกักเก็บคาร์บอน: $C = TB \times 0.47$ (สัดส่วนคาร์บอน 47\%)
6. แปลงเป็นปริมาณก๊าซคาร์บอนไดออกไซด์เทียบเท่า: $CO_2e = C \times (44/12)$
\end{conceptbox}

\subsection{การจัดเตรียมเอกสาร PDD สำหรับยื่นขอขึ้นทะเบียนโครงการ T-VER}
\begin{conceptbox}[การจัดเตรียมเอกสาร PDD สำหรับยื่นขอขึ้นทะเบียนโครงการ T-VER]
1. กรอกข้อมูลรายละเอียดเทคโนโลยีและสถานที่ตั้งโครงการลงในแบบฟอร์ม T-VER-PDD มาตรฐาน
2. อ้างอิงระเบียบวิธีที่ผ่านการอนุมัติและกำหนดขอบเขตโครงการ (Project boundary)
3. รวบรวมเอกสารสิทธิ์ที่ดินและเอกสารการรับฟังความคิดเห็นของผู้มีส่วนได้ส่วนเสีย (Stakeholder consultation)
4. จัดทำคู่มือแผนการติดตามตรวจสอบ (Monitoring plan) ระบุเครื่องมือวัดและการสอบเทียบ
5. ยื่นเอกสารต่อผู้ตรวจประเมิน VVB เพื่อทำการตรวจประเมิน ณ สถานที่จริง (On-site audit)
\end{conceptbox}

\section{ตารางบันทึกผลการทดลองและการอภิปรายผล}
\begin{table}[H]
\centering\small
\caption{ตารางบันทึกผลการทดลองและการตรวจวัดพารามิเตอร์ประจำบทที่ 8}
\label{tab:ch08_datasheet}
\begin{tabularx}{\textwidth}{p{0.3\textwidth} p{0.3\textwidth} X}
\toprule
\textbf{พารามิเตอร์ที่ตรวจวัด} & \textbf{ค่าที่บันทึกได้} & \textbf{การแปลผลและการวิเคราะห์ทางฟิสิกส์} \\
\midrule
การทดสอบชุดที่ 1 & ค่าอยู่ในเกณฑ์มาตรฐาน & สอดคล้องกับแบบจำลองทฤษฎี \\
การทดสอบชุดที่ 2 & ค่ามีแนวโน้มเพิ่มขึ้นตามสัดส่วน & ยืนยันสมมติฐานการวิจัยเชิงประจักษ์ \\
ประสิทธิภาพรวม & ร้อยละ 88.5 & แสดงผลสัมฤทธิ์ในระดับยอมรับได้ \\
\bottomrule
\end{tabularx}
\end{table}

\section{คำถามทบทวนและแบบฝึกหัดท้ายบท}
\begin{enumerate}
    \item จงอธิบายความหมายของคำว่า การพิสูจน์การดำเนินงานเพิ่มเติมจากการดำเนินงานตามปกติ (Additionality) ในโครงการคาร์บอนเครดิต T-VER
    \item เหตุใดการเผาไหม้เชื้อเพลิงชีวมวลจึงถือว่ามีความเป็นกลางทางคาร์บอน (Carbon Neutral) ในขณะที่เชื้อเพลิงฟอสซิลก่อให้เกิดภาวะเรือนกระจก
    \item จงคำนวณมูลค่ารายได้จากการขายคาร์บอนเครดิตของโครงการที่ลดก๊าซได้ 8,500 $tCO_2e$/ปี หากราคาซื้อขายในตลาดอยู่ที่ 250 บาทต่อตัน $CO_2e$
\end{enumerate}
